\documentclass{situationist-preprint}

\title{The Levelling-Down Objection}
\msid{POL-08-b}
\msfamily{POL-08}
\msdiscipline{Political Philosophy}
\mstype{Companion manuscript}
\mskeywords{levelling down; telic egalitarianism; person-affecting value; Temkin; fitting attitudes; sufficiency}

\begin{document}
\maketitle

\begin{abstract}
The levelling-down objection holds that if equality were good in itself, making the better off worse off, with no gain to anyone, would be better in one respect. Temkin replies that the objection rests on a person-affecting restriction on value that egalitarians should reject. We argue that the objection does not need that restriction. On the fitting-attitude test of POL-08-a, values that benefit no one, such as beauty, pass, and the equality produced by levelling down fails. The objection is decisive, and what it leaves standing is the sufficiency view, together with a difference principle that condemns every levelling down.
\end{abstract}

\section{Temkin's reply}

\citet{Parfit1997} put the objection thus: if inequality is bad in itself, then removing it by making the better off worse off, while making no one better off, is in one respect a change for the better. Suppose those who can see are blinded so that they become equal to those who cannot. Telic egalitarians (Definition~2.1 of POL-08-a) must say that this is better in one respect.

\citet[ch.~9]{Temkin1993} replies that the objection draws its force from the \emph{slogan} that an outcome cannot be better in any respect if it is better for no one. The slogan is false: it would also condemn values that egalitarians share with their opponents, such as the value of desert or of the survival of a species, and it has paradoxical implications in cases involving future persons. Without the slogan, he argues, the claim that levelling down is better in one respect is merely what egalitarianism says, and the egalitarian can hold it worse all things considered.

\section{The objection without the slogan}

POL-08-a does not rest on the slogan. It rests on the weak half of the fitting-attitude account (Principle~P1 of POL-08-a): an outcome is good in a respect only if some pro-attitude toward it in that respect is fitting. That principle does not restrict value to what is good for someone.

\begin{proposition}\label{p:test}
The fitting-attitude test is passed by some values that are good for no one, and failed by the equality produced by a levelling down.
\end{proposition}

\begin{proof}
A painting may be beautiful though its beauty benefits no one, as when it hangs unseen; it is fitting to be glad that it exists and to regret its destruction. So impersonal value passes the test. In a levelling down, by Lemma~3.1 of POL-08-a, the respect in which the outcome would be better is the reduced inequality, and the only change that reduces it is a loss to the better off that benefits no one. To welcome that loss as such, or to be glad of it in that respect, is the attitude of envy or spite, not a fitting response; so the test fails, as Theorem~3.2 of POL-08-a concludes.
\end{proof}

Proposition~\ref{p:test} answers Temkin's dilemma. The egalitarian cannot defend the claim that levelling down is better in one respect by pointing to impersonal values in general, since impersonal values pass a test that this one fails. And the claim that it is \enquote{worse all things considered} does not help, for the objection was never that levelling down is best overall but that there is something in it to welcome, and there is not.

\section{What the objection leaves standing}

The objection is decisive against telic egalitarianism, not against concern for the worse off. The sufficiency thesis (Definition~2.3 of POL-08-a) is untouched, since levelling down never brings anyone closer to having enough (Proposition~3.3 of POL-08-a). Nor does the objection touch the modified difference principle of POL-01, which is not egalitarian: it ranks every levelling down below the status quo (Proposition~4.1 of POL-08-a), and its priority for the worst off rests on the guarantee of the sufficiency floor, not on the badness of the gap between the worst off and others.

\section{Conclusion}

The levelling-down objection does not depend on a person-affecting restriction on value. It depends only on the claim that what is good in a respect is fit to be welcomed in that respect, and the equality produced by levelling down is not. What matters is that everyone has enough.

\begin{thebibliography}{9}
\bibitem[Parfit(1997)]{Parfit1997} Parfit, D. (1997). Equality and priority. \emph{Ratio}, 10(3), 202--221.
\bibitem[Temkin(1993)]{Temkin1993} Temkin, L.~S. (1993). \emph{Inequality}. Oxford University Press.
\end{thebibliography}

\end{document}
